%% Formule de l'operateur : Comble_interpolin
%% Dans ce fichier, les �ventuels accents sont cod�s
%% en ISO-8859-15 (aka latin9)
%% La ligne suivante donne trois exemples 
%% e accent aigu : �, e accent grave : �, e-dans-l'o : �

\vspace{1cm}
L'algorithme de comblement des valeurs manquantes fonctionne de la
mani�re suivante : \\


\begin{itemize}

\item Calcul de la moyenne glissante de $X$ sur $Nbpt$ points. 
\item Entre 2 valeurs valides de $X$, les valeurs manquantes sont combl�es
par une interpolation lin�aire de la moyenne glissante de $X$
\item Si la s�rie $X$ d�bute ou se termine par des valeurs manquantes, respectivement
la 1ere ou la derni�re valeur valide est utilis�e pour combler $X_f$
\end{itemize}
\vspace{1cm}
Le param�tre $Flag$ vaut 0 si l'�chantillon  $X_f = X $ (pas de comblement), et vaut 1
si l'�chantillon de $X_f$ a �t� combl�. 

